\documentclass[10pt,a4paper]{amsart}
\usepackage{amsmath,amssymb,amsthm,graphicx}
% Retained minimal preamble for the Stacks Project chapter flat.tex (More on Flatness).
% Verbatim excerpt of the chapter preamble: Xy-pic setup (native preamble lines 14-18) and the native
% theorem declarations with the mathematical macros (native preamble lines 154-217). Copied mechanically;
% nothing added, removed or redefined. The wrapper class and the amsmath/amssymb/amsthm/graphicx base
% packages are supplied by the assembler, so the class/externaldocument part of the native preamble is omitted.
% For commutative diagrams we use Xy-pic
\usepackage[all]{xy}

% We use 2cell for 2-commutative diagrams.
\xyoption{2cell}
% Theorem environments.
%
\theoremstyle{plain}
\newtheorem{theorem}[subsection]{Theorem}
\newtheorem{proposition}[subsection]{Proposition}
\newtheorem{lemma}[subsection]{Lemma}

\theoremstyle{definition}
\newtheorem{definition}[subsection]{Definition}
\newtheorem{example}[subsection]{Example}
\newtheorem{exercise}[subsection]{Exercise}
\newtheorem{situation}[subsection]{Situation}

\theoremstyle{remark}
\newtheorem{remark}[subsection]{Remark}
\newtheorem{remarks}[subsection]{Remarks}

\numberwithin{equation}{subsection}

% Macros
%
\def\lim{\mathop{\mathrm{lim}}\nolimits}
\def\colim{\mathop{\mathrm{colim}}\nolimits}
\def\Spec{\mathop{\mathrm{Spec}}}
\def\Hom{\mathop{\mathrm{Hom}}\nolimits}
\def\Ext{\mathop{\mathrm{Ext}}\nolimits}
\def\SheafHom{\mathop{\mathcal{H}\!\mathit{om}}\nolimits}
\def\SheafExt{\mathop{\mathcal{E}\!\mathit{xt}}\nolimits}
\def\Sch{\mathit{Sch}}
\def\Mor{\mathop{\mathrm{Mor}}\nolimits}
\def\Ob{\mathop{\mathrm{Ob}}\nolimits}
\def\Sh{\mathop{\mathit{Sh}}\nolimits}
\def\NL{\mathop{N\!L}\nolimits}
\def\CH{\mathop{\mathrm{CH}}\nolimits}
\def\proetale{{pro\text{-}\acute{e}tale}}
\def\etale{{\acute{e}tale}}
\def\QCoh{\mathit{QCoh}}
\def\Ker{\mathop{\mathrm{Ker}}}
\def\Im{\mathop{\mathrm{Im}}}
\def\Coker{\mathop{\mathrm{Coker}}}
\def\Coim{\mathop{\mathrm{Coim}}}

% Boxtimes
%
\DeclareMathSymbol{\boxtimes}{\mathbin}{AMSa}{"02}

%
% Macros for moduli stacks/spaces
%
\def\QCohstack{\mathcal{QC}\!\mathit{oh}}
\def\Cohstack{\mathcal{C}\!\mathit{oh}}
\def\Spacesstack{\mathcal{S}\!\mathit{paces}}
\def\Quotfunctor{\mathrm{Quot}}
\def\Hilbfunctor{\mathrm{Hilb}}
\def\Curvesstack{\mathcal{C}\!\mathit{urves}}
\def\Polarizedstack{\mathcal{P}\!\mathit{olarized}}
\def\Complexesstack{\mathcal{C}\!\mathit{omplexes}}
% \Pic is the operator that assigns to X its picard group, usage \Pic(X)
% \Picardstack_{X/B} denotes the Picard stack of X over B
% \Picardfunctor_{X/B} denotes the Picard functor of X over B
\def\Pic{\mathop{\mathrm{Pic}}\nolimits}
\def\Picardstack{\mathcal{P}\!\mathit{ic}}
\def\Picardfunctor{\mathrm{Pic}}
\def\Deformationcategory{\mathcal{D}\!\mathit{ef}}

% --- appended 2026-10-12 (Chief): the project own definitions of the reference/slogan/history
% display environments, copied mechanically from the shared Stacks preamble lines 8-12
% (\usepackage{verbatim} plus the three \newenvironment lines, which the chapter preamble keeps
% verbatim from the project preamble). Nothing added, removed or redefined: these environments
% carry no mathematical content and the project itself typesets them as comments. Without these
% three lines a counted statement whose own text uses them fails with Environment slogan undefined.
\usepackage{verbatim}
\newenvironment{reference}{\comment}{\endcomment}
\newenvironment{slogan}{\comment}{\endcomment}
\newenvironment{history}{\comment}{\endcomment}
\usepackage{cleveref}
\crefname{originaldefinition}{original definition}{original definitions}
\Crefname{originaldefinition}{Original definition}{Original definitions}
\makeatletter
\newlabel{definition-shrink-complete}{{05HL}{source}{}{sourcepointer-1}{}}
\newlabel{definition-shrink-complete@cref}{{[originaldefinition][0][]05HL}{[1][0][]source}{}{}{}}
\makeatother
\crefname{originalreference}{original reference}{original references}
\Crefname{originalreference}{Original reference}{Original references}
\makeatletter
\newlabel{item-shrink-on-X}{{05HD}{source}{}{sourcepointer-2}{}}
\newlabel{item-shrink-on-X@cref}{{[originalreference][0][]05HD}{[1][0][]source}{}{}{}}
\makeatother
\crefname{originalreference}{original reference}{original references}
\Crefname{originalreference}{Original reference}{Original references}
\makeatletter
\newlabel{item-shrink-on-Y}{{05HC}{source}{}{sourcepointer-3}{}}
\newlabel{item-shrink-on-Y@cref}{{[originalreference][0][]05HC}{[1][0][]source}{}{}{}}
\makeatother
\crefname{originallemma}{original lemma}{original lemmas}
\Crefname{originallemma}{Original lemma}{Original lemmas}
\makeatletter
\newlabel{lemma-base-change-complete-at-x}{{05HK}{source}{}{sourcepointer-4}{}}
\newlabel{lemma-base-change-complete-at-x@cref}{{[originallemma][0][]05HK}{[1][0][]source}{}{}{}}
\makeatother
\crefname{originallemma}{original lemma}{original lemmas}
\Crefname{originallemma}{Original lemma}{Original lemmas}
\makeatletter
\newlabel{lemma-shrink}{{05HA}{source}{}{sourcepointer-5}{}}
\newlabel{lemma-shrink@cref}{{[originallemma][0][]05HA}{[1][0][]source}{}{}{}}
\makeatother
\title{Shrinking a complete devissage: standard open neighbourhoods of the base, of the last scheme and of the point}
\author{The Stacks Project Authors}
\date{stacks-project:flat.tex:local-snapshot-20261011; GNU FDL 1.2 or later}
\begin{document}
\maketitle
\noindent\textit{Editorial note.} All delivered mathematical text is the author's own, copied from the pinned Stacks Project chapter file flat.tex; no statement or proof step was rewritten, reordered or removed.\par\medskip
\noindent\textit{Editorial note.} Preamble provenance: the chapter's own retained minimal preamble (sources/stacks-flat/retained-preamble.tex, lines 1-85) is a verbatim excerpt of the native Xy-pic setup and of the native theorem declarations and mathematical macros, followed by the project's own reference/slogan/history display-environment definitions copied mechanically from the shared preamble; the wrapper class and the amsmath/amssymb/amsthm/graphicx base packages are supplied by the assembler, and extra\_packages is left empty because everything needed is already in the retained preamble. No macro was invented or redefined.\par\medskip
\noindent\textit{Editorial note.} Same-chapter pointers are located at their \textbackslash{}label lines inside the pinned chapter file; cross-chapter pointers are located in the project's own tag list staged in the snapshot (sources/stacks-chow/tags.txt), at the verbatim lines carrying their canonical tags.\par\medskip
\noindent\textit{Allowed standard prerequisites (editorial).}\begin{itemize}
\item The notion of a standard shrinking of a complete devissage. Same chapter flat.tex, definition labelled definition-shrink-complete (statement, proof line 1363), together with the item labels item-shrink-base-complete, item-shrink-on-Y-complete and item-shrink-on-X-complete declared inside the counted statement. Fixes what has to be produced in each of the three statements: a standard shrinking whose base lies in a given standard open neighbourhood, or whose last scheme is a given W times the base, or whose source lies in a given open U.
\item Base change of a complete devissage at a point. Same chapter flat.tex, label lemma-base-change-complete-at-x, invoked at proof line 1384. Supplies, in part (1), that the base changed data still form a complete devissage at the image point, so that restricting to a standard open neighbourhood of the base gives a standard shrinking.
\item The one step shrinking lemmas. Same chapter flat.tex, label lemma-shrink, invoked at proof lines 1384 and 1389 with the item references item-shrink-on-Y and item-shrink-on-X, and again at proof line 1417. Provides one step of shrinking at the level of the last scheme or at the level of the source, which the proof iterates along the devissage in parts (2) and (3); the item labels are navigation pointers outside the counted span and carry the label string itself as their number because they are item labels, not indexed tags.
\end{itemize}\medskip
\ifdefined\hypertarget\hypertarget{sourcepointer-1}{}\fi
\noindent\textit{Source reference.} Original definition 05HL: sources/stacks-flat/pinned-original-source.tex, source lines 1341--1341. Navigation pointer for \textbackslash{}ref\{definition-shrink-complete\} used inside the counted span (line 1363); the argument that consumes it is retained in the span.\par\medskip
\ifdefined\hypertarget\hypertarget{sourcepointer-2}{}\fi
\noindent\textit{Source reference.} Original reference 05HD: sources/stacks-flat/pinned-original-source.tex, source lines 985--985. Navigation pointer for \textbackslash{}ref\{item-shrink-on-X\} used inside the counted span (line 1407); the argument that consumes it is retained in the span.\par\medskip
\ifdefined\hypertarget\hypertarget{sourcepointer-3}{}\fi
\noindent\textit{Source reference.} Original reference 05HC: sources/stacks-flat/pinned-original-source.tex, source lines 981--981. Navigation pointer for \textbackslash{}ref\{item-shrink-on-Y\} used inside the counted span (line 1389); the argument that consumes it is retained in the span.\par\medskip
\ifdefined\hypertarget\hypertarget{sourcepointer-4}{}\fi
\noindent\textit{Source reference.} Original lemma 05HK: sources/stacks-flat/pinned-original-source.tex, source lines 1313--1313. Navigation pointer for \textbackslash{}ref\{lemma-base-change-complete-at-x\} used inside the counted span (line 1384); the argument that consumes it is retained in the span.\par\medskip
\ifdefined\hypertarget\hypertarget{sourcepointer-5}{}\fi
\noindent\textit{Source reference.} Original lemma 05HA: sources/stacks-flat/pinned-original-source.tex, source lines 970--970. Navigation pointer for \textbackslash{}ref\{lemma-shrink\} used inside the counted span (line 1385); the argument that consumes it is retained in the span.\par\medskip
\begin{lemma}
\label{lemma-shrink-complete}
With assumption and notation as in
Definition \ref{definition-shrink-complete}
we have:
\begin{enumerate}
\item
\label{item-shrink-base-complete}
If $S' \subset S$ is a standard open neighbourhood of $s$, then
setting $X' = X_{S'}$, $Z'_k = Z_{S'}$ and $Y'_k = Y_{S'}$ we obtain a
standard shrinking.
\item
\label{item-shrink-on-Y-complete}
Let $W \subset Y_n$ be a standard open neighbourhood of $y$.
Then there exists a standard shrinking with $Y'_n = W \times_S S'$.
\item
\label{item-shrink-on-X-complete}
Let $U \subset X$ be an open neighbourhood of $x$.
Then there exists a standard shrinking with $X' \subset U$.
\end{enumerate}
\end{lemma}

\begin{proof}
Part (1) is immediate from
Lemmas \ref{lemma-base-change-complete-at-x} and
\ref{lemma-shrink}.

\medskip\noindent
Proof of (2). For convenience denote $X = Y_0$. We apply
Lemma \ref{lemma-shrink} (\ref{item-shrink-on-Y})
to find a standard shrinking
$S', Y'_{n - 1}, Z'_n, Y'_n$
of the one step d\'evissage of $\Coker(\alpha_{n - 1})/Y_{n - 1}/S$
at $y_{n - 1}$ with $Y'_n = W \times_S S'$. We may repeat this procedure
and find a standard shrinking
$S'', Y''_{n - 2}, Z''_{n - 1}, Y''_{n - 1}$
of the one step d\'evissage of $\Coker(\alpha_{n - 2})/Y_{n - 2}/S$
at $y_{n - 2}$ with $Y''_{n - 1} = Y'_{n - 1} \times_S S''$.
We may continue in this manner until we obtain
$S^{(n)}, Y^{(n)}_0, Z^{(n)}_1, Y^{(n)}_1$.
At this point it is clear that we obtain our desired standard shrinking
by taking $S^{(n)}$, $X^{(n)}$, $Z_k^{(n - k)} \times_S S^{(n)}$, and
$Y_k^{(n - k)} \times_S S^{(n)}$ with the desired property.

\medskip\noindent
Proof of (3). We use induction on the length of the complete
d\'evissage. First we apply
Lemma \ref{lemma-shrink} (\ref{item-shrink-on-X})
to find a standard shrinking
$S', X', Z'_1, Y'_1$
of the one step d\'evissage of $\mathcal{F}/X/S$ at $x$
with $X' \subset U$. If $n = 1$, then we are done.
If $n > 1$, then by induction we can find a standard shrinking
$S''$, $Y''_1$, $Z''_k$, and $Y''_k$ of the complete d\'evissage
$(Z_k, Y_k, i_k, \pi_k, \mathcal{G}_k, \alpha_k, z_k, y_k)_{k = 2, \ldots, n}$
of $\Coker(\alpha_1)/Y_1/S$ at $x$ such that
$Y''_1 \subset Y'_1$. Using
Lemma \ref{lemma-shrink} (\ref{item-shrink-on-Y})
we can find $S''' \subset S'$, $X''' \subset X'$, $Z'''_1$ and
$Y'''_1 = Y''_1 \times_S S'''$ which is a standard shrinking.
The solution to our problem is to take
$$
S''', X''', Z'''_1, Y'''_1, Z''_2 \times_S S''',
Y''_2 \times_S S''', \ldots, Z''_n \times_S S''', Y''_n \times_S S'''
$$
This ends the proof of the lemma.
\end{proof}

\end{document}
